% utilities
\usepackage{etoolbox}
\usepackage{calc}

% bibliography
\usepackage[T1]{fontenc}
\usepackage[utf8]{inputenc}
\usepackage[backend=biber, style=authoryear, maxcitenames=2, url=true, doi=false, isbn=false, backref=true]{biblatex}
\addbibresource{references.bib}

% Custom inline citation format: [Author, Title, Year] in gray, clickable
\DeclareCiteCommand{\cite}
  {\textcolor{gray}{[}\usebibmacro{prenote}}
  {\bibhyperref{\textcolor{gray}{\printnames{labelname}%
   \setunit{\addcomma\addspace}%
   \mkbibemph{\printfield{title}}%
   \setunit{\addcomma\addspace}%
   \printfield{year}}}}
  {\multicitedelim}
  {\usebibmacro{postnote}\textcolor{gray}{]}}

% packages
\usepackage[boxed, noline, boxruled, noend]{algorithm2e}
%\RestyleAlgo{ruled}
\SetKwInput{KwInput}{Input}
\SetKwInput{KwOutput}{Output}
\DontPrintSemicolon
\usepackage{tcolorbox}
\usepackage{listings}
\usepackage{amsmath}
\usepackage{amsthm}
\usepackage{amsfonts}
\usepackage{amssymb}
\usepackage{amsxtra}
\usepackage{bbm}
\usepackage{mathrsfs}
\usepackage{verbatim}
\usepackage{subfig}
\usepackage{ragged2e}
\usepackage{array}
\usepackage{hhline}
\usepackage{cancel}
\usepackage{units}
\usepackage{color}
\usepackage{colortbl}
\usepackage{bm}
\usepackage{animate}
\usepackage[absolute]{textpos}
\usepackage{multirow}
\usepackage{multicol}
\usepackage{booktabs}
\usepackage{hyperref}
\usepackage{tabularx}
\usepackage{xcolor}
\definecolor{links}{HTML}{2A1B81}
\hypersetup{colorlinks,linkcolor=,urlcolor=links}
\usepackage[font=small,justification=centering]{caption}
\usepackage{enumerate}       
\usepackage{enumitem}
\usepackage[normalem]{ulem}
\usepackage{pgfplots}
\pgfplotsset{compat=1.18}
\usepackage{pgfplots}
\usepackage{tikz}
\usetikzlibrary{fit}
\usetikzlibrary{backgrounds}
\usetikzlibrary{shapes}
\usetikzlibrary{calc}
\usetikzlibrary{arrows}
\usetikzlibrary{arrows.meta}
\usetikzlibrary{positioning}
\usetikzlibrary{decorations.pathreplacing}
\usetikzlibrary{patterns}
\usepackage{ifthen}
%\usepackage{subfig}
\usepackage{subcaption}
\usepackage{numerica}
%\usepackage{qbordermatrix}
\usepackage{scalerel}
\usepackage{stackengine}
\usepackage{eurosym}

% Commands
\newcommand\blfootnote[1]{%
  \begingroup
  \renewcommand\thefootnote{}\footnote{#1}%
  \addtocounter{footnote}{-1}%
  \endgroup
}

% centered box (ignoring margins)
\newcommand\cbox[1]{%
\makebox[\textwidth][c]{%
#1%
}}

%% vertically centered
\newcommand{\vcentered}[1]{\begingroup\setbox0=\hbox{#1}\parbox{\wd0}{\box0}\endgroup}

% Math short commands
%\newcommand{\E}[1]{\mathbb{E}\left[#1\right]} %Expected value
%\newcommand\gauss[2]{1/(#2*sqrt(2*pi))*exp(-((x-#1)^2)/(2*#2^2))} % Gauss function, parameters mu and sigma
%\newcommand\gauss[2]{1/(#2*sqrt(2*pi))*exp(-((x-#1)^2)/(2*#2^2))} % Gauss function, parameters mu and sigma
\pgfmathdeclarefunction{gauss}{2}{%
  \pgfmathparse{1/(#2*sqrt(2*pi))*exp(-((x-#1)^2)/(2*#2^2))}%
}

%\newcommand{\E}[1][]{%
%\ifthenelse{\equal{#1}{}}{\mathbb{E}}{\mathbb{P}\left[{#1}\right]}%
%}
\newcommand{\infdiv}{D_\textrm{KL}\infdivx}
\newcommand{\E}{\mathbb{E}} 
\newcommand{\R}{\mathbb{R}} 
\newcommand{\N}{\mathbb{N}} 
\newcommand{\Z}{\mathbb{Z}} 
%\newcommand{\Var}[1]{\mathrm{Var}\left[#1\right]} %Variance value
%\newcommand{\Cov}[1]{\mathrm{Cov}\left[#1\right]} %Variance value
\newcommand{\El}[2]{\mathbb{E}_{#2}\left[#1\right]} %Expected value with lower index
%\newcommand{\Pb}[1]{\mathbb{P}\left[{#1}\right]} %Probability
%Probability
\newcommand{\Scope}[1][]{%
\ifthenelse{\equal{#1}{}}{\text{Scope}}{\text{Scope}\left[{#1}\right]}%
}
\newcommand{\Pb}[1][]{%
\ifthenelse{\equal{#1}{}}{\mathbb{P}}{\mathbb{P}\left[{#1}\right]}%
}
\newcommand{\Pa}[2]{\text{Pa}^{#1}_{#2}} %Parents
\newcommand{\Ch}[2]{\text{Ch}^{#1}_{#2}} %Children
\newcommand{\Desc}[2][]{%
\ifthenelse{\equal{#2}{}}{\text{Desc}^{#1}}{\text{Desc}^{#1}_{#2}}%
}
\newcommand{\Anc}{\text{Anc}}
\newcommand{\Mor}{\mathcal{M}}
\newcommand{\Nb}{\text{Nb}}
\newcommand{\ND}[2]{\text{NonDesc}^{#1}_{#2}} %Non-descendants
\newcommand{\maxmarg}{\text{MaxMarg}}
%\newcommand{\T}{^{\mkern-1.5mu\mathsf{T}}} %Transpose
\DeclareMathOperator*{\argmax}{arg\,max} % arg max
\DeclareMathOperator*{\argmin}{arg\,min} % arg min
\newcommand{\abs}[1]{|#1|}
\newcommand{\norm}[1]{||#1||}
\newcommand{\MB}{\text{MB}}
\newcommand{\indep}{\bot}
\newcommand{\Indep}{\mathcal{I}}
\newcommand{\la}{\langle}
\newcommand{\ra}{\rangle}
\newcommand{\T}{\top}

\newcommand\tr[1]{#1^{\textrm{T}}}
\newcommand\normn[1]{\lVert #1 \rVert}
\newcommand\fnorm[1]{\left\lVert #1 \right\rVert_\text{F}}
\newcommand\fnormn[1]{\lVert #1 \rVert_\text{F}}
\newcommand\rsim{\stackrel{\mathrm{row}}{\sim}}
\newcommand\csim{\stackrel{\mathrm{col}}{\sim}}


\newcommand{\idx}{\text{idx}}

\definecolor{emph}{rgb}{0.31,0.5,0.74}
\newcommand\emphdef[1]{\textbf{\textcolor{emph}{#1}}}
\newcommand\emphcall[2]{\if*#1\textcolor{emph}{\textsc{#2}}\else\textcolor{emph}{\textsc{#1}(#2)}\fi} % define method call (* w/o args)
\ifthenelse{\equal{show only notes}{\notes}}
  {\newcommand{\putat}[3]{}}
  {%
    \newcommand{\putat}[3]{%
      \begin{textblock*}{0mm}[0,0](#1,#2)%
        #3
      \end{textblock*}%
    }%
  }
\newcommand\sourceref[1]{\putat{.1em}{9.4cm}{\mbox{{\tiny\textcolor{gray}{#1}}}}}

% tikz styles
\tikzset{ rand_var/.style = {rectangle, rounded corners, draw, align=center} }
\tikzset{ prior_var/.style = {ellipse, draw, align=center, minimum width=0.7cm, minimum height=0.5cm, inner sep=0pt} }
\tikzset{ factor_var/.style = {rectangle, draw} }
\tikzset{ every edge/.style = {very thick} }


% ----------------
% VECTOR VARIABLES
% ----------------
\newcommand\vect[1]{{\boldsymbol{#1}}}
\newcommand\va{\vect{a}}
\newcommand\vb{\vect{b}}
\newcommand\vc{\vect{c}}
\newcommand\vd{\vect{d}}
\newcommand\ve{\vect{e}}
\newcommand\vf{\vect{f}}
\newcommand\vg{\vect{g}}
\newcommand\vh{\vect{h}}
\newcommand\vi{\vect{i}}
\newcommand\vj{\vect{j}}
\newcommand\vk{\vect{k}}
\newcommand\vl{\vect{l}}
\newcommand\vm{\vect{m}}
\newcommand\vn{\vect{n}}
\newcommand\vo{\vect{o}}
\newcommand\vp{\vect{p}}
\newcommand\vq{\vect{q}}
\newcommand\vr{\vect{r}}
\newcommand\vs{\vect{s}}
\newcommand\vt{\vect{t}}
\newcommand\vu{\vect{u}}
\newcommand\vv{\vect{v}}
\newcommand\vw{\vect{w}}
\newcommand\vx{\vect{x}}
\newcommand\vy{\vect{y}}
\newcommand\vz{\vect{z}}
\newcommand\vzero{\vect{0}}
\newcommand\vone{\vect{1}}

\newcommand\valpha{\vect{\alpha}}
\newcommand\vbeta{\vect{\beta}}
\newcommand\veps{\vect{\epsilon}}
\newcommand\vdelta{\vect{\delta}}
\newcommand\veta{\vect{\eta}}
\newcommand\vtau{\vect{\tau}}
\newcommand\vtheta{\vect{\theta}}
\newcommand\vmu{\vect{\mu}}
\newcommand\vsigma{\vect{\sigma}}
\newcommand\vpi{\vect{\pi}}
\newcommand\vlambda{\vect{\lambda}}

% ----------------
% MATRIX VARIABLES
% ----------------
\newcommand\mA{\vect{A}}
\newcommand\mB{\vect{B}}
\newcommand\mC{\vect{C}}
\newcommand\mD{\vect{D}}
\newcommand\mE{\vect{E}}
\newcommand\mF{\vect{F}}
\newcommand\mG{\vect{G}}
\newcommand\mH{\vect{H}}
\newcommand\mI{\vect{I}}
\newcommand\mJ{\vect{J}}
\newcommand\mK{\vect{K}}
\newcommand\mL{\vect{L}}
\newcommand\mM{\vect{M}}
\newcommand\mN{\vect{N}}
\newcommand\mO{\vect{O}}
\newcommand\mP{\vect{P}}
\newcommand\mQ{\vect{Q}}
\newcommand\mR{\vect{R}}
\newcommand\mS{\vect{S}}
\newcommand\mT{\vect{T}}
\newcommand\mU{\vect{U}}
\newcommand\mV{\vect{V}}
\newcommand\mW{\vect{W}}
\newcommand\mX{\vect{X}}
\newcommand\mY{\vect{Y}}
\newcommand\mZ{\vect{Z}}
\newcommand\mzero{\vect{0}}

\newcommand{\mEps}{{\ensuremath{\vect{\Epsilon}}}}
\newcommand{\mSigma}{{\ensuremath{\vect{\Sigma}}}}
\newcommand{\mLambda}{{\ensuremath{\vect{\Lambda}}}}
\newcommand{\mPsi}{{\ensuremath{\vect{\Psi}}}}

% ------------------
% BLACKBOARD LETTERS
% ------------------
\newcommand\bA{\mathbb{A}}
\newcommand\bB{\mathbb{B}}
\newcommand\bC{\mathbb{C}} % set of complex numbers
\newcommand\bD{\mathbb{D}}
\newcommand\bE{\mathbb{E}}
\newcommand\bF{\mathbb{F}}
\newcommand\bG{\mathbb{G}}
\newcommand\bH{\mathbb{H}}
\newcommand\bI{\mathbb{I}}
\newcommand\bJ{\mathbb{J}}
\newcommand\bK{\mathbb{K}}
\newcommand\bL{\mathbb{L}}
\newcommand\bM{\mathbb{M}}
\newcommand\bN{\mathbb{N}} % set of positive integers
\newcommand\bO{\mathbb{O}}
\newcommand\bP{\mathbb{P}}
\newcommand\bQ{\mathbb{Q}} % set of rationals
\newcommand\bR{\mathbb{R}} % set of reals
\newcommand\bS{\mathbb{S}}
\newcommand\bT{\mathbb{T}}
\newcommand\bU{\mathbb{U}}
\newcommand\bV{\mathbb{V}}
\newcommand\bW{\mathbb{W}}
\newcommand\bX{\mathbb{X}}
\newcommand\bY{\mathbb{Y}}
\newcommand\bZ{\mathbb{Z}} % set of all integers

% -------------
% FANCY LETTERS
% -------------

% requires mathrsfs package
\newcommand\xA{\mathscr{A}}
\newcommand\xB{\mathscr{B}}
\newcommand\xC{\mathscr{C}}
\newcommand\xD{\mathscr{D}}
\newcommand\xE{\mathscr{E}}
\newcommand\xF{\mathscr{F}}
\newcommand\xG{\mathscr{G}}
\newcommand\xH{\mathscr{H}}
\newcommand\xI{\mathscr{I}}
\newcommand\xJ{\mathscr{J}}
\newcommand\xK{\mathscr{K}}
\newcommand\xL{\mathscr{L}}
\newcommand\xM{\mathscr{M}}
\newcommand\xN{\mathscr{N}}
\newcommand\xO{\mathscr{O}}
\newcommand\xP{\mathscr{P}}
\newcommand\xQ{\mathscr{Q}}
\newcommand\xR{\mathscr{R}}
\newcommand\xS{\mathscr{S}}
\newcommand\xT{\mathscr{T}}
\newcommand\xU{\mathscr{U}}
\newcommand\xV{\mathscr{V}}
\newcommand\xW{\mathscr{W}}
\newcommand\xX{\mathscr{X}}
\newcommand\xY{\mathscr{Y}}
\newcommand\xZ{\mathscr{Z}}

% --------------------
% CALLIGRAPHIC LETTERS
% --------------------
\DeclareMathAlphabet{\mathcal}{OMS}{cmsy}{m}{n}
\newcommand\cA{\mathcal{A}}
\newcommand\cB{\mathcal{B}}
\newcommand\cC{\mathcal{C}}
\newcommand\cD{\mathcal{D}}
\newcommand\cE{\mathcal{E}}
\newcommand\cF{\mathcal{F}}
\newcommand\cG{\mathcal{G}}
\newcommand\cH{\mathcal{H}}
\newcommand\cI{\mathcal{I}}
\newcommand\cJ{\mathcal{J}}
\newcommand\cK{\mathcal{K}}
\newcommand\cL{\mathcal{L}}
\newcommand\cM{\mathcal{M}}
\newcommand\cN{\mathcal{N}}
\newcommand\cO{\mathcal{O}}
\newcommand\cP{\mathcal{P}}
\newcommand\cQ{\mathcal{Q}}
\newcommand\cR{\mathcal{R}}
\newcommand\cS{\mathcal{S}}
\newcommand\cT{\mathcal{T}}
\newcommand\cU{\mathcal{U}}
\newcommand\cV{\mathcal{V}}
\newcommand\cW{\mathcal{W}}
\newcommand\cX{\mathcal{X}}
\newcommand\cY{\mathcal{Y}}
\newcommand\cZ{\mathcal{Z}}

% -------------------------
% BOLD CALLIGRAPHIC LETTERS
% -------------------------
\DeclareMathAlphabet\mathbfcal{OMS}{cmsy}{b}{n}
\newcommand\tA{\mathbfcal{A}}
\newcommand\tB{\mathbfcal{B}}
\newcommand\tC{\mathbfcal{C}}
\newcommand\tD{\mathbfcal{D}}
\newcommand\tE{\mathbfcal{E}}
\newcommand\tF{\mathbfcal{F}}
\newcommand\tG{\mathbfcal{G}}
\newcommand\tH{\mathbfcal{H}}
\newcommand\tI{\mathbfcal{I}}
\newcommand\tJ{\mathbfcal{J}}
\newcommand\tK{\mathbfcal{K}}
\newcommand\tL{\mathbfcal{L}}
\newcommand\tM{\mathbfcal{M}}
\newcommand\tN{\mathbfcal{N}}
\newcommand\tO{\mathbfcal{O}}
\newcommand\tP{\mathbfcal{P}}
\newcommand\tQ{\mathbfcal{Q}}
\newcommand\tR{\mathbfcal{R}}
\newcommand\tS{\mathbfcal{S}}
\newcommand\tT{\mathbfcal{T}}
\newcommand\tU{\mathbfcal{U}}
\newcommand\tV{\mathbfcal{V}}
\newcommand\tW{\mathbfcal{W}}
\newcommand\tX{\mathbfcal{X}}
\newcommand\tY{\mathbfcal{Y}}
\newcommand\tZ{\mathbfcal{Z}}



% -----------------
% LETTERS WITH HATS
% -----------------

% use amsxtra package!

\newcommand\Ahat{{\hat A}}
\newcommand\Bhat{{\hat B}}
\newcommand\Chat{{\hat C}}
\newcommand\Dhat{{\hat D}}
\newcommand\Ehat{{\hat E}}
\newcommand\Fhat{{\hat F}}
\newcommand\Ghat{{\hat G}}
\newcommand\Hhat{{\hat H}}
\newcommand\Ihat{{\hat I}}
\newcommand\Jhat{{\hat J}}
\newcommand\Khat{{\hat K}}
\newcommand\Lhat{{\hat L}}
\newcommand\Mhat{{\hat M}}
\newcommand\Nhat{{\hat N}}
\newcommand\Ohat{{\hat O}}
\newcommand\Phat{{\hat P}}
\newcommand\Qhat{{\hat Q}}
\newcommand\Rhat{{\hat R}}
\newcommand\Shat{{\hat S}}
\newcommand\That{{\hat T}}
\newcommand\Uhat{{\hat U}}
\newcommand\Vhat{{\hat V}}
\newcommand\What{{\hat W}}
\newcommand\Xhat{{\hat X}}
\newcommand\Yhat{{\hat Y}}
\newcommand\Zhat{{\hat Z}}

\newcommand\ahat{{\hat a}}
\newcommand\bhat{{\hat b}}
\newcommand\chat{{\hat c}}
\newcommand\dhat{{\hat d}}
\newcommand\ehat{{\hat e}}
\newcommand\fhat{{\hat f}}
\newcommand\ghat{{\hat g}}
\newcommand\hhat{{\hat h}}
\newcommand\ihat{{\hat i}}
\newcommand\jhat{{\hat j}}
\newcommand\khat{{\hat k}}
\newcommand\lhat{{\hat l}}
\newcommand\mhat{{\hat m}}
\newcommand\nhat{{\hat n}}
\newcommand\ohat{{\hat o}}
\newcommand\phat{{\hat p}}
\newcommand\qhat{{\hat q}}
\newcommand\rhat{{\hat r}}
\newcommand\shat{{\hat s}}
\newcommand\that{{\hat t}}
\newcommand\uhat{{\hat u}}
\newcommand\vhat{{\hat v}}
\newcommand\what{{\hat w}}
\newcommand\xhat{{\hat x}}
\newcommand\yhat{{\hat y}}
\newcommand\zhat{{\hat z}}

\newcommand\rhohat{{\hat\rho}}

% -----------------
% LETTERS WITH BARS
% -----------------

\accentedsymbol\Abar{{\bar A}}
\accentedsymbol\Bbar{{\bar B}}
\accentedsymbol\Cbar{{\bar C}}
\accentedsymbol\Dbar{{\bar D}}
\accentedsymbol\Ebar{{\bar E}}
\accentedsymbol\Fbar{{\bar F}}
\accentedsymbol\Gbar{{\bar G}}
\accentedsymbol\Hbar{{\bar H}}
\accentedsymbol\Ibar{{\bar I}}
\accentedsymbol\Jbar{{\bar J}}
\accentedsymbol\Kbar{{\bar K}}
\accentedsymbol\Lbar{{\bar L}}
\accentedsymbol\Mbar{{\bar M}}
\accentedsymbol\Nbar{{\bar N}}
\accentedsymbol\Obar{{\bar O}}
\accentedsymbol\Pbar{{\bar P}}
\accentedsymbol\Qbar{{\bar Q}}
\accentedsymbol\Rbar{{\bar R}}
\accentedsymbol\Sbar{{\bar S}}
\accentedsymbol\Tbar{{\bar T}}
\accentedsymbol\Ubar{{\bar U}}
\accentedsymbol\Vbar{{\bar V}}
\accentedsymbol\Wbar{{\bar W}}
\accentedsymbol\Xbar{{\bar X}}
\accentedsymbol\Ybar{{\bar Y}}
\accentedsymbol\Zbar{{\bar Z}}

\accentedsymbol\abar{{\bar a}}
\accentedsymbol\bbar{{\bar b}}
\accentedsymbol\cbar{{\bar c}}
\accentedsymbol\dbar{{\bar d}}
\accentedsymbol\ebar{{\bar e}}
\accentedsymbol\fbar{{\bar f}}
\accentedsymbol\gbar{{\bar g}}
\makeatletter
\@ifundefined{hbar}{}{
        \let\oldhbar\hbar % we are going to overwrite an existing command; keep the old version
        \let\hbar\@undefined
}
\makeatother
\accentedsymbol\hbar{{\bar h}}
\accentedsymbol\ibar{{\bar i}}
\accentedsymbol\jbar{{\bar j}}
\accentedsymbol\kbar{{\bar k}}
\accentedsymbol\lbar{{\bar l}}
\accentedsymbol\mbar{{\bar m}}
\accentedsymbol\nbar{{\bar n}}
\makeatletter
\@ifundefined{obar}{}{
        \let\oldobar\obar % we are going to overwrite an existing command; keep the old version
        \let\obar\@undefined
}
\makeatother
\accentedsymbol{\obar}{{\bar o}}
\accentedsymbol\pbar{{\bar p}}
\accentedsymbol\qbar{{\bar q}}
\accentedsymbol\rbar{{\bar r}}
\accentedsymbol\sbar{{\bar s}}
\accentedsymbol\tbar{{\bar t}}
\accentedsymbol\ubar{{\bar u}}
\accentedsymbol\vbar{{\bar v}}
\accentedsymbol\wbar{{\bar w}}
\accentedsymbol\xbar{{\bar x}}
\accentedsymbol\ybar{{\bar y}}
\accentedsymbol\zbar{{\bar z}}

% -------------
% GREEK LETTERS
% -------------
\newcommand\eps{\epsilon}


% ---------------------------------
% MATRIX/VECTOR VARIABLES WITH HATS
% ---------------------------------

\accentedsymbol\mAhat{{\hat\mA}}
\accentedsymbol\mBhat{{\hat\mB}}
\accentedsymbol\mChat{{\hat\mC}}
\accentedsymbol\mDhat{{\hat\mD}}
\accentedsymbol\mEhat{{\hat\mE}}
\accentedsymbol\mFhat{{\hat\mF}}
\accentedsymbol\mGhat{{\hat\mG}}
\accentedsymbol\mHhat{{\hat\mH}}
\accentedsymbol\mIhat{{\hat\mI}}
\accentedsymbol\mJhat{{\hat\mJ}}
\accentedsymbol\mKhat{{\hat\mK}}
\accentedsymbol\mLhat{{\hat\mL}}
\accentedsymbol\mMhat{{\hat\mM}}
\accentedsymbol\mNhat{{\hat\mN}}
\accentedsymbol\mOhat{{\hat\mO}}
\accentedsymbol\mPhat{{\hat\mP}}
\accentedsymbol\mQhat{{\hat\mQ}}
\accentedsymbol\mRhat{{\hat\mR}}
\accentedsymbol\mShat{{\hat\mS}}
\accentedsymbol\mThat{{\hat\mT}}
\accentedsymbol\mUhat{{\hat\mU}}
\accentedsymbol\mVhat{{\hat\mV}}
\accentedsymbol\mWhat{{\hat\mW}}
\accentedsymbol\mXhat{{\hat\mX}}
\accentedsymbol\mYhat{{\hat\mY}}
\accentedsymbol\mZhat{{\hat\mZ}}

\accentedsymbol\vahat{{\hat\va}}
\accentedsymbol\vbhat{{\hat\vb}}
\accentedsymbol\vchat{{\hat\vc}}
\accentedsymbol\vdhat{{\hat\vd}}
\accentedsymbol\vehat{{\hat\ve}}
\accentedsymbol\vfhat{{\hat\vf}}
\accentedsymbol\vghat{{\hat\vg}}
\accentedsymbol\vhhat{{\hat\vh}}
\accentedsymbol\vihat{{\hat\vi}}
\accentedsymbol\vjhat{{\hat\vj}}
\accentedsymbol\vkhat{{\hat\vk}}
\accentedsymbol\vlhat{{\hat\vl}}
\accentedsymbol\vmhat{{\hat\vm}}
\accentedsymbol\vnhat{{\hat\vn}}
\accentedsymbol\vohat{{\hat\vo}}
\accentedsymbol\vphat{{\hat\vp}}
\accentedsymbol\vqhat{{\hat\vq}}
\accentedsymbol\vrhat{{\hat\vr}}
\accentedsymbol\vshat{{\hat\vs}}
\accentedsymbol\vthat{{\hat\vt}}
\accentedsymbol\vuhat{{\hat\vu}}
\accentedsymbol\vvhat{{\hat\vv}}
\accentedsymbol\vwhat{{\hat\vw}}
\accentedsymbol\vxhat{{\hat\vx}}
\accentedsymbol\vyhat{{\hat\vy}}
\accentedsymbol\vzhat{{\hat\vz}}

\urlstyle{sf}
\let\oldhref\href
\renewcommand\href[2]{\oldhref{#1}{\textcolor{emph}{\uline{#2}}}} % emphasize href
\let\oldurl\url
\renewcommand\url[1]{\href{#1}{#1}} % emphasize url
\newcommand\email[1]{\href{mailto:#1}{\texttt{#1}}} % email address

\newcommand{\largesep}{\setlength\itemsep{2ex}}
\newcommand{\smallsep}{\setlength\itemsep{1ex}}

\renewcommand\:{\colon} % for use with \sset, etc.

% pause after math
\newcommand\mpause{\vskip-\baselineskip\pause}

\newcommand\fit[1]{\mbox{#1\hspace{-5cm}}} % fit into single line

\newenvironment<>{hblock}[1]{%
  \begin{actionenv}#2%
    \def\insertblocktitle{#1}%
    \par%
    \mode<presentation>{%
      \setbeamercolor{heading}{fg=black,bg=white}
    }%
    \setbeamertemplate{blocks}[rounded][shadow=false]
    \setbeamercolor{block title}{fg=structure,bg=white}
    \setbeamercolor{block body}{fg=black,bg=white}
    \usebeamertemplate{block begin}}
  {\par%
    \usebeamertemplate{block end}%
  \end{actionenv}}

 \newcommand\corresponds{\mathrel{\stackon[1.5pt]{=}{\stretchto{%
 \scalerel*[\widthof{=}]{\wedge}{\rule{1ex}{3ex}}}{0.5ex}}}}

% yellow "observe" box
\newcommand\observe[2][]{\begin{tikzpicture} \node[fill=yellow,draw=black,#1] {#2} ; \end{tikzpicture}}

\newcommand\raitem[1]{\begin{tabular}{@{}p{.5cm}@{}p{\linewidth-.5cm}@{}}
  \ra & #1
\end{tabular}}

\newlength\tikzdrawcx%
\newlength\tikzdrawcy%
\newcommand\tikzdrawc[5]{%
  \setlength\tikzdrawcx{#2-#4/2}%
  \setlength\tikzdrawcy{#3-#5/2}%
  \putat{\tikzdrawcx}{\tikzdrawcy}{%
    \begin{tikzpicture}#1\end{tikzpicture}
  }%
}

% centered ellipse [stroke width], color,  centerx, centery, width, height
\newcommand\ellipsec[6][.7pt]{%
  \tikzdrawc{\draw [draw=#2,line width=#1] ($(#5/2,#6/2)$) ellipse (#5/2 and #6/2) ;}{#3}{#4}{#5}{#6}%
}

\newcommand{\dd}{\; \mathrm{d}} % correct spacing and formating


% don't skip before math environment
\newcommand\nomathskip{\abovedisplayskip=0pt \abovedisplayshortskip=0pt}


% ------
% MACROS
% ------

\newcommand\dom{\operatorname{dom}}
\newcommand\Dom{\operatorname{Dom}}
\newcommand\Val{\operatorname{Val}}
\newcommand\Ber{\operatorname{Ber}}
\newcommand\Mu{\operatorname{Mu}}
\newcommand\Dir{\operatorname{Dir}}
\newcommand\Beta{\operatorname{Beta}}
\newcommand\Cat{\operatorname{Cat}}
\newcommand\Bin{\operatorname{Bin}}
\newcommand\sort{\alpha}
\newcommand\tuple[1]{\langle#1 \rangle}
\newcommand\seq[1]{\langle#1 \rangle}
\newcommand\Wit{\operatorname{Wit}}
\newcommand\MWit{\operatorname{MWit}}
\newcommand\Why{\operatorname{Why}}
\newcommand\MWhy{\operatorname{MWhy}}
\newcommand\Min{\operatorname{Min}}
\newcommand\odds{\operatorname{odds}}
\newcommand\PosBool{\operatorname{PosBool}}
\newcommand\supp{\operatorname{supp}}
\newcommand\setc[1]{#1^\mathrm{c}}
\newcommand\sep{\operatorname{sep}}
\newcommand\cnf{\operatorname{cnf}}
\newcommand\ip[1]{\langle #1 \rangle}
\newcommand\logit{\operatorname{logit}}
\let\oldodot\odot
\let\odot\undefined
\DeclareMathOperator{\odot}{\oldodot}

% --------------
% EDITING MACROS
% --------------
\newcommand{\todo}[1]{\textcolor{red}{#1}} % todos
\newcommand{\co}[1]{\textcolor{blue}{#1}} % comments
\newcommand{\eat}[1]{} % TO MAKE LARGE BLOCKS OF TEXT INVISIBLE

\newcommand{\red}[1]{\textcolor{red}{#1}}
\newcommand{\blue}[1]{\textcolor{blue}{#1}}
\newcommand{\green}[1]{\textcolor{green}{#1}}
\newcommand{\dgreen}[1]{\textcolor{green!70!black}{#1}}

% ----
% BIAS
% ----
\newcommand{\xbias}[2]{\operatorname{bias}\td{#1}{#2}} % Bias^?_?
\newcommand{\bias}[1]{\xbias**\left[\macrospace{}#1\macrospace{}\right]} % Bias[?], automatic brackets
\newcommand{\biasn}[1]{\xbias**[\macrospace{}#1\macrospace{}]} % Bias[?], normal brackets

% ---------------------
% SUPER- AND SUBSCRIPTS
% ---------------------
\newcommand{\td}[2]{\if*#1\else^{#1}\fi\if*#2\else_{#2}\fi} % ^?_?, *=ignore

% -------------
% SPACING
% -------------
\newcommand\macrospace{}

% ----------
% EQUALITIES
% ----------
\newcommand{\eqdef}{\stackrel{\mathrm{def}}{=}} % definining equality
\newcommand{\eqdot}{\doteq} % APPROACHES THE LIMIT (asymptotic equality)
\newcommand{\eqdist}{\equiv} % equality in distribution


% --------
% VARIANCE
% --------
\newcommand{\xvar}[2]{\operatorname{var}\td{#1}{#2}} % Var^?_?
\newcommand{\var}[1]{\xvar**\left[\macrospace{}#1\macrospace{}\right]} % Var[?], automatic brackets
\newcommand{\varn}[1]{\xvar**[\macrospace{}#1\macrospace{}]} % Var[?], normal brackets
\newcommand{\varbig}[1]{\xvar**\bigl[\macrospace{}#1\macrospace{}\bigr]} % Var[?], big brackets
\newcommand{\varBig}[1]{\xvar**\Bigl[\macrospace{}#1\macrospace{}\Bigr]} % Var[?], Big brackets

\newcommand{\xhatvar}[2]{\operatorname{\hat{v}ar}\td{#1}{#2}} % ~Var^?_?
\newcommand{\hatvar}[1]{\xhatvar**\left[\macrospace{}#1\macrospace{}\right]} % ~Var[?], automatic brackets
\newcommand{\hatvarn}[1]{\xhatvar**[\macrospace{}#1\macrospace{}]} % ~Var[?], normal brackets
\newcommand{\hatvarBig}[1]{\xhatvar**\Bigl[\macrospace{}#1\macrospace{}\Bigr]} % ~Var[?], Big brackets

\newcommand{\varf}[3]{\xvar{#1}{#2}{\left[\macrospace{}#3\macrospace{}\right]}} % Var^?_?[?], automatic brackets
\newcommand{\varfn}[3]{\xvar{#1}{#2}{[\macrospace{}#3\macrospace{}]}} % E^?_?[?], normal brackets
\newcommand{\varfbig}[3]{\xvar{#1}{#2}{\bigl[\macrospace{}#3\macrospace{}\bigr]}} % E^?_?[?], big brackets
\newcommand{\varfBig}[3]{\xvar{#1}{#2}{\Bigl[\macrospace{}#3\macrospace{}\Bigr]}} % E^?_?[?], Big brackets
\newcommand{\varfbigg}[3]{\xvar{#1}{#2}{\biggl[\macrospace{}#3\macrospace{}\biggr]}} % E^?_?[?], bigg

% ----------
% COVARIANCE
% ----------
\newcommand{\xcov}[2]{\operatorname{Cov}\td{#1}{#2}} % Cov^?_?
\newcommand{\cov}[1]{\xcov**\left[\macrospace{}#1\macrospace{}\right]} % Cov[?], automatic brackets
\newcommand{\covn}[1]{\xcov**[\macrospace{}#1\macrospace{}]} % Cov[?], normal brackets
\newcommand{\covBig}[1]{\xcov**\Bigl[\macrospace{}#1\macrospace{}\Bigr]} % Cov[?], Big brackets

\newcommand{\xhatcov}[2]{\operatorname{\hat{C}ov}\td{#1}{#2}} % ~Cov^?_?
\newcommand{\hatcov}[1]{\xhatcov**\left[\macrospace{}#1\macrospace{}\right]} % ~Cov[?], automatic brackets
\newcommand{\hatcovn}[1]{\xhatcov**[\macrospace{}#1\macrospace{}]} % ~Cov[?], normal brackets
\newcommand{\hatcovBig}[1]{\xhatcov**\Bigl[\macrospace{}#1\macrospace{}\Bigr]} % ~Cov[?], Big brackets

% --------------
% STANDARD ERROR
% --------------

\newcommand{\xse}[2]{\operatorname{SE}\td{#1}{#2}} % SE^?_?
\newcommand{\se}[1]{\xse**\left[\macrospace{}#1\macrospace{}\right]} % SE[?], automatic brackets
\newcommand{\sen}[1]{\xse**[\macrospace{}#1\macrospace{}]} % SE[?], normal brackets

\newcommand{\xhatse}[2]{\operatorname{\hat{S}E}\td{#1}{#2}} % ~SE^?_?
\newcommand{\hatse}[1]{\xhatse**\left[\macrospace{}#1\macrospace{}\right]} % ~SE[?], automatic brackets
\newcommand{\hatsen}[1]{\xhatse**[\macrospace{}#1\macrospace{}]} % ~SE[?], normal brackets

% --------------
% STANDARD ERROR
% --------------

\newcommand{\xsd}[2]{\operatorname{SD}\td{#1}{#2}} % SD^?_?
\newcommand{\sd}[1]{\xsd**\left[\macrospace{}#1\macrospace{}\right]} % SD[?], automatic brackets
\newcommand{\sdn}[1]{\xsd**[\macrospace{}#1\macrospace{}]} % SD[?], normal brackets

\newcommand{\xhatsd}[2]{\operatorname{\hat{S}D}\td{#1}{#2}} % ~SD^?_?
\newcommand{\hatsd}[1]{\xhatsd**\left[\macrospace{}#1\macrospace{}\right]} % ~SD[?], automatic brackets
\newcommand{\hatsed}[1]{\xhatsd**[\macrospace{}#1\macrospace{}]} % ~SD[?], normal brackets


% ------------------------
% COEFFICIENT OF VARIATION
% ------------------------
\newcommand{\xcv}[2]{\operatorname{CV}\td{#1}{#2}} % Cv^?_?
\newcommand{\cv}[1]{\xcv**\left[\macrospace{}#1\macrospace{}\right]} % Cv[?], automatic brackets
\newcommand{\cvn}[1]{\xcv**[\macrospace{}#1\macrospace{}]} % Var[?], normal brackets
\newcommand{\cvbig}[1]{\xcv**\bigl[\macrospace{}#1\macrospace{}\bigr]} % Cv[?], big brackets
\newcommand{\cvBig}[1]{\xcv**\Bigl[\macrospace{}#1\macrospace{}\Bigr]} % Cv[?], Big brackets

% -----------------
% MEAN SQUARE ERROR
% -----------------
\newcommand{\xmse}[2]{\operatorname{mse}\td{#1}{#2}} % MSE^?_?
\newcommand{\mse}[1]{\xmse**\left[\macrospace{}#1\macrospace{}\right]} % MSE[?], automatic brackets
\newcommand{\msen}[1]{\xmse**[\macrospace{}#1\macrospace{}]} % MSE[?], normal brackets

% ----------------------
% ROOT MEAN SQUARE ERROR
% ----------------------
\newcommand{\xrmse}[2]{\operatorname{rmse}\td{#1}{#2}} % RMSE^?_?
\newcommand{\rmse}[1]{\xrmse**\left[\macrospace{}#1\macrospace{}\right]} % RMSE[?], automatic brackets
\newcommand{\rmsen}[1]{\xrmse**[\macrospace{}#1\macrospace{}]} % RMSE[?], normal brackets

% ----------------------
% AVERAGE RELATIVE ERROR
% ----------------------
\newcommand{\xare}[2]{\operatorname{ARE}\td{#1}{#2}} % ARE^?_?
\newcommand{\are}[1]{\xare**\left[\macrospace{}#1\macrospace{}\right]} % ARE[?], automatic brackets
\newcommand{\aren}[1]{\xare**[\macrospace{}#1\macrospace{}]} % ARE[?], normal brackets

% ---------------------
% BINOMIAL DISTRIBUTION
% ---------------------

% B(k; N,q)
\newcommand{\pbinom}[3]{B\!\left(#1;\,#2,#3\right)} % binomial probability, automatic brackets

% ---------------------------
% HYPERGEOMETRIC DISTRIBUTION
% ---------------------------
\newcommand{\phyper}[4]{H\!\left(#1;\,#2,#3,#4\right)} % hypergeometric probability, automatic brackets


% -------------
% CARDINALITIES
% -------------
\newcommand{\sz}[1]{\lvert#1\rvert}   % cardinality of set
\newcommand{\card}[1]{\lvert#1\rvert} % cardinality of set
\newcommand{\Sz}[1]{\lVert#1\rVert}   % cardinality of set, distinct (not used in thesis!)
\newcommand{\Card}[1]{\lVert#1\rVert} % cardinality of set, distinct (not used in thesis!)

% -----------------
% MATRIX OPERATIONS
% -----------------
\newcommand\xdiag{\operatorname{diag}}
\newcommand\diag[1]{\xdiag\left(#1\right)}    % diagonal matrix
\newcommand\diagn[1]{\xdiag(#1)}    % diagonal matrix, normal brackets
\newcommand\xrank{\operatorname{rank}}
\newcommand\rank[1]{\xrank\left(#1\right)}
\newcommand\rankn[1]{\xrank(#1)}

% ------
% TABLES
% ------

\newcolumntype{Y}{>{\centering\arraybackslash}X}

% forces an image to the top of a table cell
\def\imagetop#1{\raisebox{-\height+\baselineskip}{#1}}

% forces an image to the center of an item
\def\imagec#1{\raisebox{-.5\height+.5\baselineskip}{#1}}

% environment for tightly formatted tables (relations); use via \relation* commands
\newenvironment{rel}{\bgroup\tabular}{\endtabular\egroup}
\newcommand\relname[2]{\multicolumn{#1}{l}{\hspace{-.3em}\parbox{0cm}{\mbox{\textcolor{LightMidnightBlue}{#2}}}}\\}
\newcommand\relhead[1]{\cellcolor{gray}\textcolor{white}{#1}}
\newcommand\reltop[1]{\hhline{*{#1}-}}
\newcommand\relmid[1]{\hhline{*{#1}-}}
\newcommand\relbot[1]{\hhline{*{#1}-}}
\newcommand\reltopl[1]{\hhline{~*{#1}-}}
\newcommand\relmidl[1]{\hhline{~*{#1}-}}
\newcommand\relbotl[1]{\hhline{~*{#1}-}}
\newcommand\reltopr[1]{\hhline{*{#1}-~}}
\newcommand\relmidr[1]{\hhline{*{#1}-~}}
\newcommand\relbotr[1]{\hhline{*{#1}-~}}
\newcommand\reltoprl[1]{\hhline{~*{#1}-~}}
\newcommand\relmidrl[1]{\hhline{~*{#1}-~}}
\newcommand\relbotrl[1]{\hhline{~*{#1}-~}}
\newcolumntype{T}{>{\columncolor{white}}c} % right math

% a relation: name (optional), number of columns, content
\newcommand\relation[3][*]{{\setlength{\tabcolsep}{2pt}\begin{rel}{|*{#2}{T|}}\if*#1\else\relname{#2}{#1}\fi\reltop{#2}#3\\\relbot{#2}\end{rel}}}

% a relation: name (optional), number of columns, content
\newcommand\relationb[3][*]{{\setlength{\tabcolsep}{2pt}\begin{rel}{|*{#2}{T|}}\if*#1\else\relname{#2}{#1}\fi\reltop{#2}#3\\\relbot{#2}\multicolumn{#2}{l}{}\end{rel}}}

% a relation with tuples on left side: name (optional), number of columns, content
\newcommand\relationl[3][*]{{\setlength{\tabcolsep}{2pt}\begin{rel}{r@{\hspace{.2em}}|*{#2}{T|}}\if*#1\else\multicolumn{1}{c}{}&\relname{#2}{\hspace{-.1em}#1}\fi\reltopl{#2}#3\\\relbotl{#2}\end{rel}}}

% a relation with tuples on right side: name (optional), number of columns, content
\newcommand\relationr[3][*]{{\setlength{\tabcolsep}{2pt}\begin{rel}{|*{#2}{T|}l}\if*#1\else\relname{#2}{\hspace{-.1em}#1}\fi\reltopr{#2}#3\\\relbotr{#2}\end{rel}}}

% a relation with tuples on left and right side: name (optional), number of columns, content
\newcommand\relationrl[3][*]{{\setlength{\tabcolsep}{2pt}\begin{rel}{r|*{#2}{T|}l}\if*#1\else\multicolumn{1}{c}{}&\relname{#2}{\hspace{-.1em}#1}\fi\reltoprl{#2}#3\\\relbotrl{#2}\end{rel}}}

% a table with a header; use via \tab* commands
\newcommand\tabhead[1]{\textbf{#1}}
\newcommand\tabtop[1]{\hhline{*{#1}-}}
\newcommand\tabmid[1]{\hhline{*{#1}-}}
\newcommand\tabbot[1]{\hhline{*{#1}-}}
\newcommand\tabtopr[1]{\hhline{*{#1}-~}}
\newcommand\tabmidr[1]{\hhline{*{#1}-~}}
\newcommand\tabbotr[1]{\hhline{*{#1}-~}}

% a table, args: name, columns, content
\newcommand\tab[3]{\begin{tabular}{#2} \tabtop{#1}#3\\ \tabbot{#1}\end{tabular}}
% a table with tuples right, args: name, columns, content
\newcommand\tabr[3]{\begin{tabular}{#2} \tabtopr{#1}#3\\ \tabbotr{#1}\end{tabular}}

% ----
% SETS
% ----
\renewcommand\:{\colon} % for use with \sset, etc.
\newcommand{\sset}[1]{\left\{\,#1\,\right\}} % { ? }, automatic brackets
\newcommand{\ssets}[1]{\left\{#1\right\}} % {?}, automatic brackets
\newcommand{\ssetn}[1]{\{\,#1\,\}} % { ? }, normal brackets

\newcommand{\comp}[1]{\bar{#1}} % complement

\newcommand{\powerset}[1]{\cP\!\left(#1\right)} % powerset, automatic brackets

\newcommand\adots{{\ensuremath ..}} % short array notation as in A[i..j]

\newcommand{\Nat}{{\ensuremath \bN}}
\newcommand{\Bool}{{\ensuremath \{0,1\}}}

% default colors for background of marked matrix rows/columns/intersects
\definecolor{rcol}{rgb}{0.5,1,0.5}
\definecolor{ccol}{rgb}{1,1,.5}
\definecolor{rccol}{rgb}{0.125,.97,.0625}
